\documentclass[instructions.tex]{subfiles}
\begin{document}
 
\chapter{Avis pour l'éducation chrétienne de la jeunesse.}
 
Il est peu d’art dont les règles soient en même temps plus difficiles, moins connues, et si malheureusement négligées dans la pratique. Il seroit à souhaiter que les pasteurs, rassemblant de temps en temps, et à part, ceux de leurs paroissiens qui se trouvent chargés d’élever des enfans, entrassent avec eux dans le détail de leurs principaux devoirs, leur indiquant la manière de les remplir, et d’éviter les défauts qui se commettent communément dans cet important objet. Ne pouvant dire ici tout ce qui seroit nécessaire, nous nous contenterons des avis suivans.
 
Former l’esprit des enfans par l’instruction; leur cœur, par les vertus chrétiennes, et les placer dans un état convenable, où ils rendront des services à la société, tout en remplissant les devoirs qu’impose la religion, voilà en gros à quoi se réduisent les devoirs des parens pour l’éducation religieuse de leurs enfans.
 
\primo Ne pas envoyer les enfans à l’école, n’avoir pas soin qu’ils apprennent le catéchisme, qu’ils assistent à l’explication qu’en donnent les pasteurs, et les laisser croupir dans une ignorance totale des vérités de la religion, c’est de la part des parens un meurtre spirituel, plus criminel que le meurtre temporel qu’ils commettraient en les étouffant immédiatement après leur avoir procuré le baptême : par ce dernier crime, ils feroient leur bonheur, en envoyant leur âme dans le sein de Dieu; par le premier, ils font leur malheur, les plaçant à coup sûr dans le chemin de l’enfer.
 
Mais tout en voulant instruire les enfans de ce qui concerne la religion, il faut descendre à leur portée, y aller par degrés, rendre la doctrine sensible par des comparaisons, autant qu’il se peut, et surtout se garder de mêler des erreurs aux vérités saintes, des superstitions aux pratiques approuvées par l’église, des exagérations aux règles douces de la morale évangélique.
 
Laissez donc aux pasteurs le soin d’exposer la doctrine des mystères, dont vous n’avez peut-être pas vous-même une idée bien juste et bien précise; ou contentez-vous de rappeler à vos enfans ce que vous êtes sûr d’avoir bien retenu des explications pastorales; à moins de cela, tenez-vous-en à inculquer littéralement ce qui est dans le catéchisme. Ceci regarde toutes les personnes chargées de l’éducation, qui n’ont pas fait les études théologiques nécessaires au catéchiste.
 
Défiez-vous de vos lumières sur les règles de la morale, craignant de donner pour conseil ce qui est de précepte, pour précepte ce qui n’est que de conseil, pour mal ce qui est indifférent de soi, et pour bien ou pour indifférent ce que la sévérité de l’Évangile réprouve. Que de mauvaises consciences remontent à l’enfance !
 
\secundo A l’égard des vices, et de ce qui est défendu, soyez très-prudent; ne tranchez point sur le véniel et le mortel; la distance entre l’un et l’autre échappe quelquefois aux docteurs les plus subtils, et la décision en cette matière est souvent très-difficile, dit saint Augustin. Dire, par exemple, que tous les menteurs sont les enfans du démon, c’est évidemment mentir soi-même, induire en erreur un enfant, lui donner une décision qui lui formera une conscience erronée, d’où résulteront des péchés graves où ils ne seroient d’ailleurs que très-légers.
 
Touchant les vertus, faites entendre à vos enfans qu’elles sont belles en elles-mêmes, agréables à Dieu, utiles à la société; des ornemens précieux qui embellissent le cœur de celui qui les pratique, et qu’elles lui obtiendront une grande gloire dans le ciel. Mais gardez-vous de tomber dans un défaut pernicieux, dont certains parens ne s’éloignent pas assez : lesquels, tout en stimulant leurs enfans au bien, en quelque genre que ce soit, au lieu d’y exciter leur émulation par des motifs purs et salutaires, les y portent par le désir de paraître, d’être loués, admirés. De pareilles leçons nourrissent et développent l’orgueil; elles font la vanité et la présomption, ce qui s’ensuit; elles accoutument les enfans à faire le bien par des motifs pharisiaques et condamnables, par conséquent à mériter des châtimens plutôt que des récompenses, même pour le bien qu’ils font, à cause de l’orgueil qui y préside et qui l’infecte devant Dieu.
 
Éloignez de la satisfaction à votre enfant, quand il a bien agi, obtenu du succès, fait une action digne de louange; mais dites-lui aussi-là, que c’est de Dieu que vient tout bien, puisque c’est lui qui nous donne les moyens, la force, en un mot, la volonté, et qu’il nous aide en tous les cas où nous faisons le bien. Il faut, par conséquent, qu’il lui rende toute la gloire de nos succès et de nos bonnes œuvres; ensuite, remerciez Dieu et engagez votre enfant à le remercier aussi, et à toujours profiter des dons de la bonté divine, afin de faire de mieux en mieux.
 
Quant à ce qui concerne la modestie, mettez votre adresse, d’abord à éloigner l’occasion des défauts contraires, sans que votre enfant s’en doute et qu’il s’en aperçoive; ensuite, quand il lui arrivera de manquer, mais sans savoir encore distinguer ce qui est bien de ce qui est mal en cette matière, dites-lui de se tenir ce qui est mal en cette matière, dites-lui de se tenir ce qui est mal, et le conduire à ce qu’il faut. Allez toujours ainsi par degré, à mesure que la raison se développera, et suivant les connoissances que votre enfant prendra, en prenant garde de ne point prévenir ses connoissances à cet égard, et de lui apprendre le mal qu’il ignore encore heureusement. Des avertissemens indiscrets sont souvent cause beaucoup de mal, et n’ont procuré aucun bien en ce point. Enfin, surveillez avec prudence, prévenez les dangers, c’est-à-dire les occasions, et imitez en ceci la police, qui ne fait des règlemens et des arrêts qu’après qu’elle a découvert des désordres.
 
\tertio A ces mesures sages joignons une surveillance exacte de l’enfance et de la jeunesse. Un grand dépôt nous est confié dans la personne de nos enfans, dit saint Chrysostome, s’adressant aux pères et mères; conservons-le avec un extrême soin, et faisons tout pour empêcher qu’un voleur, plein d’astuce et de finesse ne nous l’enlève\footnote{Hom. 5, in cap. 2. Ep. 1. ad Tim. — (2) 1. Tim. 5. 8.}. Saint Augustin compare les familles à l’Église, et il dit que les parens en sont comme les évêques, parce qu’ils doivent y exercer une surveillance d’autorité et de sollicitude, à peu près comme les premiers pasteurs dans leurs diocèses. Saint Paul assure que celui qui n’a pas soin des siens, et principalement de sa maison, a renoncé à la foi, et est pire qu’un infidèle\footnote{1. Tim. 5. 8.}. Surveillez donc pour détourner le mal, dès qu’il se montre; en écarter le danger, dès qu’il menace; appliquer les enfans à la pratique des vertus de leur âge, les empêcher d’en contracter les défauts et les vices; surveillance pour la prière, la fréquentation de l’école, du catéchisme, des offices, des sacremens, la tenue respectueuse dans le lieu saint; surveillance pour éloigner les compagnies mauvaises, les discours dangereux, les manières grossières, les actions scandaleuses, les livres et les brochures capables de pervertir les mœurs, de faire perdre la foi; surveillance continuelle, qui les suive partout, qui ne se repose point, qui invente mille moyens de ne pas les perdre de vue. La police a ses agens pour ne pas s’exposer à des surprises; ayons-nous aussi des personnes discrètes et sages pour les seconder et les avertir. Enfin, il faut que vos enfans apprennent, par expérience, que votre œil est constamment ouvert sur eux, soit que vous soyez près, soit que vous vous trouviez éloignés, et qu’ils ignorent de qui vous viennent les renseignemens que vous acquérez par le témoignage d’autrui.
 
Ceux qui négligent de surveiller leurs enfans, et qui leur laissent une liberté trop grande, répondront devant Dieu des vices et des maux incalculables qui auront trouvé leur source dans ces deux déplorables excès.
 
Mais, puisque nous parlons de liberté, accordez-en assez à vos enfans, en votre présence, pour qu’ils vous laissent apercevoir le fond de leur âme, pour que leur esprit se forme, leur cœur se dilate, leurs organes se fortifient, leur santé s’affermisse. N’exigez donc point qu’ils soient devant vous tranquilles comme des vieillards, immobiles comme des termes, taciturnes comme des statues, permisez-leur même un peu de langue. Il leur faut de l’exercice, procurez-leur des jeux innocens, du plaisir, faites qu’ils en trouvent chez vous, autant et même de plus délicieux qu’ailleurs, afin qu’ils se plaisent à la maison, qu’ils y reviennent avec empressement, quand il n’y survient pas, et qu’ils s’y tiennent volontiers, lorsque rien ne les appelle au-dehors.
 
Quant aux veillées, aux assemblées et aux promenades, réglez vos enfans sur ce qui en a été dit ci-devant (chap. III, pag. 137 et suiv.). Ce que j’en dirai ici est essentiel que tout se passe sous vos yeux ou en présence d’autres personnes dignes de votre confiance, et qui vous remplaceront. Quels malheurs et quels désordres n’a point occasionés dans tous les temps le défaut de cette mesure si rigoureusement nécessaire !
 
Ne produisez point vos enfans de trop bonne heure dans le monde : ce sont des fleurs tendres qu’un souffle légèrement empoisonné flétrit, et fait aisément périr. Vouloir qu’ils soient de toutes les parties de plaisir, les conduire aux bals, aux spectacles, c’est entreprendre de renverser l’édifice qu’on s’est efforcé de bâtir en eux, empêcher le succès de leurs études, en remplissant leur cœur d’affections étrangères, et les faire renouer avec le démon, aux pompes duquel ils ont renoncé dans leur baptême. Lisez aussi à cet égard les chapitres II et III de ce livre.
 
Cependant, tout en accordant une liberté convenable à vos enfans en votre présence, ne leur y permettez pas la licence. Qu’ils se rappellent toujours le respect qu’ils vous doivent : ne leur souffrez donc à votre égard, ni à l’égard de leurs autres supérieurs, aucun manquement, aucune parole arrogante, bien moins les injures, des menaces. Qu’ils ne vous disent, ni à ceux qu’ils doivent respecter, toi; cette expression, autorisée dans un temps de délire, où tout conspire contre l’ordre, les mœurs et les principes consacrateurs de la société, doit être bannie de la bouche des enfans; on voit assez pour quoi.
 
Cependant ce seroit en vain qu’on exercerait la surveillance la plus soignée et la plus exacte, si on n’ajoutait pas à ce moyen si nécessaire la correction faite à propos et d’une manière prudente. La correction est donc encore un devoir indispensablement attaché aux fonctions de l’éducation. Écoutons là-dessus saint-Esprit dans les livres sacrés. « Celui qui épargne la verge à son fils, le hait\footnote{Prov. 13. 24.}; Celui qui l’aime, y a souvent recours\footnote{Eccl. 30. 1.}. Courbez son cou dans la jeunesse et frappez ses côtés pendant qu’il est enfant, de peur qu’il ne s’endurcisse, ne vous écoute plus, et ne devienne pour votre ame un sujet de douleur\footnote{Ibid. 29. 17.}. Vous le frapperez avec la verge, et vous délivrerez son âme de l’enfer\footnote{Prov. 23. 14.}. La verge et les menaces donnent la sagesse; mais un enfant abandonné à sa volonté fera la confusion de sa mère\footnote{Prov. 29. 14.}. »
 
Remplissez donc ce devoir important, mais remplissez-le avec prudence.
 
Pour cela, commencez par parler raison à votre enfant, lui montrant clairement les motifs qui doivent le détourner de la faute qu’il a commise, du dédégoût qu’il contracte. Répétez-lui plusieurs fois, dans l’occasion, encore avec douceur, les mêmes avis; parlez-lui ensuite avec plus de force, menacez-le, s’il ne se corrige pas; mais menacez-le de punitions graduelles, commençant par les plus légères, et tenez parole, s’il retombe : cette fidélité à votre parole, soit que vous menaciez, soit que vous fassiez des promesses. N’en venez aux coups qu’à l’extrémité, et au moins après avoir employé des moyens plus doux, à moins après avoir averti votre enfant, n’ayant une faute de nature à mériter de suite un châtiment rigoureux.
 
Mais qu’un air d’intérêt et de douceur préside toujours à vos avis, à vos avertissemens, à vos menaces, et aux autres corrections que vous emploierez. Si sur-tout, vous vous trouvez obligé d’en venir à des peines ou à des châtimens graves, que votre enfant sente que c’est de votre part à regret, et parce que les moyens plus doux ont été appliqués inutilement.
 
Soyez juste; ne punissez pas pour une bagatelle, pour un défaut léger d’attention, pour un vase fragile brisé par inadvertance. Conservez, autant que possible, en ce point comme en tout autre, une grande impartialité envers tous vos enfans. La prédilection engendre les jalousies, les haines, les divisions, on en a vu dans l’histoire sacrée et ailleurs de bien funestes exemples.
 
On voit des parens qui ne parlent jamais avec bonté à leurs enfans, qui s’emportent, prennent feu, éclatent en menaces, en propos grossiers, et frappent à tort et à droit, pour des riens comme pour des choses importantes. Une conduite si déraisonnable gâte tout; elle rend l’enfant pusillanime; elle lui inspire une humeur être jamais, et qui servira dont il ne se guérira peut-être jamais, et qui sera pour lui une source de fautes; elle le révolte, aliène son cœur et lui porte à souhaiter la mort de ses parens, auprès desquels se croît toujours en présence de l’ennemi.
 
Sur toute chose, ne vous laissez jamais aller à des malédictions : Dieu les entend, et quelquefois il les exécute d’une manière aussi prompte que terrible.
 
Une veuve, chargée d’une famille de dix enfans, ayant été grièvement offensée par l’aîné, sans que les autres s’y opposassent, alla se prosterner devant Dieu, dans le baptistère de l’église, et lui demanda que ses enfans, obligés de quitter leur patrie et d’errer dans des terres étrangères, devinssent pour le genre humain un exemple de l’erreur et d’horreur. Aussitôt, tandis qu’un tremblement horrible dans tous ses membres, les autres éprouvent successivement, selon l’ordre de leur naissance, et dans l’année, le même sort; ils abandonnent la maison paternelle, vaguent au loin sans se fixer nulle part, épouvantent le monde par le spectacle qu’ils présentent dans leurs personnes, et par le récit déplorable de leur malheur. Désespérée d’avoir été ainsi exaucée, ne tenant pas contre les remords de sa conscience, et accablée d’opprobre, la mère se pendit. Deux de ses enfans furent guéris miraculeusement à Hippone, en présence de saint Augustin, qui raconte le fait plus au long\footnote{T. 6. serm. 322. et 323.}. Mais qu’est-il besoin de remonter si haut dans le dernier siècle, une mère travaillant avec sa fille au bord d’une rivière de France. Un essaim de mouches l’importune, et lui fatiguant tout-à-coup l’autre : la fille, comme possédée, noyée à l’instant, renversée dans la rivière et entraînée, à quelque distance de là, dans les roues d’un moulin, qui la froissent d’une manière horrible. La mère, tout éplorée, alla sur-le-champ se jeter aux pieds d’un ministre de la réconciliation, pour obtenir miséricorde dans le sacrement de pénitence.
 
V. Au reste, vous ferez peu, ou plutôt vous ne ferez rien, malgré tous vos efforts, si vos enfans aperçoivent de la contradiction entre vos avis, entre ce que vous exigez d’eux, et votre conduite. Ils sont portés à l’imitation, parce qu’ils ne trouvent pas en eux-mêmes assez de ressource\footnote{Ibid. 29. 17.}. Si donc vos exemples ne concordent pas avec vos leçons et vos commandemens, ils laisseront ceux-ci, dès qu’ils le pourront, pour suivre ceux-là, qu’ils trouveront aussi plus commodes, plus conformes aux inclinations de la nature, et moins gênans. D’ailleurs, et voici le raisonnement qui ne manque pas de faire de très-bonne heure : Si ce que mes parens me prescrivent étoit avantageux et nécessaire, ils le seroit également pour eux; donc, la religion, les préceptes divins et ecclésiastiques, les obligations également : puisqu’ils s’en dispensent eux-mêmes, ils regardent donc toutes ces pratiques, ces devoirs et ces vertus, comme inutiles et mal fondés. Or, qui oblige, moi, pour en savoir et en faire plus qu’eux ?
 
VI. A l’égard du choix d’un état pour vos enfans, prenez garde d’étendre trop loin vos droits, de vous opposer à des inclinations honnêtes et venues du ciel, de consulter la cupidité, l’ambition, de vaines espérances, plutôt que la volonté de Dieu et la véritable vocation de vos enfans. Une résistance longue et déraisonnable à leur établissement a quelquefois jeté de jeunes garçons et de jeunes filles dans le désespoir, et dans une suite aussi affreuse que prolongée de crimes et de désordres. Mais pouvez-vous consentir que vos enfans s’associent avec des personnes d’une autre religion que la religion catholique ? Gardez-vous-en bien, et ne négligez rien pour dissoudre les nœuds qui tendroient à de semblables unions, dès que vous connoitrez qu’il s’en forme de cette nature dans votre famille. Quelle société peut-il y avoir entre la lumière et les ténèbres, et entre Jésus-Christ et Bélial? dit saint Paul\footnote{2. Cor. 6. 14.}. Il est impossible, dit le grand saint Hilaire; et la raison elle-même ne le permet pas, que les choses qui se répugnent aïent entre elles de la convenance; que celles qui sont opposées se réunissent, que la vérité s’allie avec le mensonge; que la lumière et la nuit subsistent ensemble dans un même lieu\footnote{Opus in Matth. 65. Chrys.}. L’Église a toujours eu en horreur des mariages où le salut des époux et celui des enfans est constamment exposé à des dangers imminens.
 
VII. Mais en vain l’homme plante et arrose, si Dieu lui-même ne donne l’accroissement. Aux soins et aux précautions sages, joignez encore la prière : c’est un devoir pour vous, et ce qui achèvera d’attirer la bénédiction du ciel sur votre famille, si votre conduite n’y met pas d’obstacle.
 
VIII. Enfin, seroit-il besoin de rappeler aux parens leur devoir d’offrir à Dieu leurs enfans, de se hâter de procurer le baptême aussitôt qu’ils sont nés, d’appeler auprès d’eux le ministre sacré quand ils sont malades, afin qu’ils reçoivent de lui la bénédiction prescrite, s’ils n’ont pas encore la raison, les sacremens de pénitence et d’extrême-onction, s’ils ont péché déjà ? L’insouciance de quelques parens à plusieurs de ces égards, a pu causer la perte éternelle à plusieurs enfans morts ou sans baptême (et la chose n’est pas douteuse), ou sans confession, à l’âge de 7, 8, 9 et même 10 ans, (et le péril est grand). O Dieu! que de parens mériteraient plutôt le nom d’assassin, que celui de père et de mère !
 
\section{Résolutions}
 
\primo Regardez l’éducation de vos enfans comme la plus noble de vos occupations et le plus important de vos devoirs. 

\secundo Commencez dès ce jour à réformer vos mauvaises manières à cet égard, et tout ce qu’il y a de répréhensible dans l’extérieur de votre conduite. 

\tertio Bannissez de votre maison jusqu’à l’ombre de la médisance. 

\prayer{O mon Dieu! vous m’avez donné des enfans, aidez-moi à les élever, dans votre amour, pour le ciel, et faites que je me sanctifie avec eux.}
 
\end{document}
